% GENERATED FILE. Edit canonical lessons, then run npm run generate:books.
% canonical-source-hash: fnv1a64:c688fb18f5ad3923
% canonical-lessons: BN-C513-kota, BN-C513-khosa, BN-C513-kora-grate, BN-C513-makha, BN-C513-kuchono

\chapter{Chopping and Kneading}
\label{ch:bn-chopping-and-kneading}
\hlchaptermodality{513}

\begin{chapteropening}
\textbf{By the end of this chapter:} \emph{I can say how food is chopped, peeled, kneaded and rolled out.}

Complete the last lesson of chapter 513: I can say how food is chopped, peeled, kneaded and rolled out.
\end{chapteropening}

\section[\texorpdfstring{\bn{কোটা} (koṭā)}{কোটা (koṭā)}]{\bn{কোটা} (koṭā) — to chop (vegetables or fish)}

\label{lesson:BN-C513-kota}



\begin{quote}
Before the new one: say the Bengali for a lid, then the Bengali for to grind.
\end{quote}

\subsection*{You'll want to know: \bn{কোটা}}
\textbf{\bn{কোটা}} — \emph{koṭā} — \textquotedblleft{}to chop (vegetables or fish)\textquotedblright{}.

\begin{quote}
\bn{আমি} \bn{রাতের} \bn{জন্য} \bn{সবজি} \bn{কুটলাম।}
\end{quote}

\begin{quote}
\emph{Āmi rāter jônyô sôbji kuṭlām.}
\end{quote}

\begin{quote}
I chopped vegetables for dinner.
\end{quote}

Past (he/she …): \bn{তিনি} \bn{মাছ} \bn{কুটলেন}, \emph{tini māchh kuṭlen}.

The stem vowel rises in the past: \bn{কুটলাম}, I chopped.

\begin{cousinweb}
\textbf{\bn{কোটা}} (\emph{koṭā}) is a tadbhava: CDIAL lists Bengali \emph{kuṭā}, \textquotedblleft{}to cut up for cooking\textquotedblright{}, and \emph{koṭā} under Sanskrit \emph{kuṭṭayati}, \textquotedblleft{}crushes, pounds\textquotedblright{} (CDIAL 3241).
\end{cousinweb}

\begin{grammarlens}[title={What you've built}]
One more word for describing what you lived and felt.
\end{grammarlens}

\subsection*{Your turn}
\begin{itemize}
\raggedright
  \item \emph{Say it:} \emph{koṭā}
  \item \emph{Say it:} \emph{koṭā}, once more
  \item \emph{Say it:} \emph{Āmi rāter jônyô sôbji kuṭlām.} — I chopped vegetables for dinner.
  \item \emph{Recall:} say the Bengali for a lid, then the Bengali for to grind, then say \emph{koṭā} again
\end{itemize}

\begin{tcolorbox}[breakable,colback=teal!4,colframe=teal!35!black,title={Before you move on}]
What does \bn{কোটা} mean? (\textquotedblleft{}To chop (vegetables or fish)\textquotedblright{}.) Say it once more.
\end{tcolorbox}
\section[\texorpdfstring{\bn{খোসা} (khosā)}{খোসা (khosā)}]{\bn{খোসা} (khosā) — a peel; a husk}

\label{lesson:BN-C513-khosa}



\begin{quote}
Before the new one: say the Bengali for to grind, then the Bengali for to chop.
\end{quote}

\subsection*{You'll want to know: \bn{খোসা}}
\textbf{\bn{খোসা}} — \emph{khosā} — \textquotedblleft{}a peel; a husk\textquotedblright{}.

\begin{quote}
\bn{আমি} \bn{আলুর} \bn{খোসা} \bn{বালতিতে} \bn{ফেললাম।}
\end{quote}

\begin{quote}
\emph{Āmi ālur khosā bāltite phellām.}
\end{quote}

\begin{quote}
I threw the potato peel into the bucket.
\end{quote}

\begin{cousinweb}
\textbf{\bn{খোসা}} (\emph{khosā}) comes from a reconstructed \emph{khosa}; CDIAL lists Bengali \emph{khosā}, \textquotedblleft{}husk\textquotedblright{} (CDIAL 3948).
\end{cousinweb}

\begin{grammarlens}[title={What you've built}]
One more word for describing what you lived and felt.
\end{grammarlens}

\subsection*{Your turn}
\begin{itemize}
\raggedright
  \item \emph{Say it:} \emph{khosā}
  \item \emph{Say it:} \emph{khosā}, once more
  \item \emph{Say it:} \emph{Āmi ālur khosā bāltite phellām.} — I threw the potato peel into the bucket.
  \item \emph{Recall:} say the Bengali for to grind, then the Bengali for to chop, then say \emph{khosā} again
\end{itemize}

\begin{tcolorbox}[breakable,colback=teal!4,colframe=teal!35!black,title={Before you move on}]
What does \bn{খোসা} mean? (\textquotedblleft{}A peel; a husk\textquotedblright{}.) Say it once more.
\end{tcolorbox}
\section[\texorpdfstring{\bn{কোরা} (korā)}{কোরা (korā)}]{\bn{কোরা} (korā) — to grate (coconut); to scrape}

\label{lesson:BN-C513-kora-grate}



\begin{quote}
Before the new one: say the Bengali for to chop, then the Bengali for a peel.
\end{quote}

\subsection*{You'll want to know: \bn{কোরা}}
\textbf{\bn{কোরা}} — \emph{korā} — \textquotedblleft{}to grate (coconut); to scrape\textquotedblright{}.

\begin{quote}
\bn{দিদিমা} \bn{পিঠের} \bn{জন্য} \bn{নারকেল} \bn{কুরলেন।}
\end{quote}

\begin{quote}
\emph{Didimā piṭher jônyô nārkel kurlen.}
\end{quote}

\begin{quote}
My grandmother grated coconut for the pithe.
\end{quote}

Past (he/she …): \bn{তিনি} \bn{নারকেল} \bn{কুরলেন}, \emph{tini nārkel kurlen}.

The stem vowel rises in the past: \bn{কুরলেন}, she grated.

\begin{cousinweb}
\textbf{\bn{কোরা}} (\emph{korā}) comes from a reconstructed \emph{kōrati}, \textquotedblleft{}scrapes\textquotedblright{}; CDIAL lists Bengali \emph{korā} and \emph{kurā}, \textquotedblleft{}to scrape, to grate\textquotedblright{} (CDIAL 3530).
\end{cousinweb}

\begin{grammarlens}[title={What you've built}]
One more word for describing what you lived and felt.
\end{grammarlens}

\subsection*{Your turn}
\begin{itemize}
\raggedright
  \item \emph{Say it:} \emph{korā}
  \item \emph{Say it:} \emph{korā}, once more
  \item \emph{Say it:} \emph{Didimā piṭher jônyô nārkel kurlen.} — My grandmother grated coconut for the pithe.
  \item \emph{Recall:} say the Bengali for to chop, then the Bengali for a peel, then say \emph{korā} again
\end{itemize}

\begin{tcolorbox}[breakable,colback=teal!4,colframe=teal!35!black,title={Before you move on}]
What does \bn{কোরা} mean? (\textquotedblleft{}To grate (coconut); to scrape\textquotedblright{}.) Say it once more.
\end{tcolorbox}
\section[\texorpdfstring{\bn{মাখা} (mākhā)}{মাখা (mākhā)}]{\bn{মাখা} (mākhā) — to knead; to mix by hand}

\label{lesson:BN-C513-makha}



\begin{quote}
Before the new one: say the Bengali for a peel, then the Bengali for to grate.
\end{quote}

\subsection*{You'll want to know: \bn{মাখা}}
\textbf{\bn{মাখা}} — \emph{mākhā} — \textquotedblleft{}to knead; to mix by hand\textquotedblright{}.

\begin{quote}
\bn{আমি} \bn{রুটির} \bn{জন্য} \bn{আটা} \bn{মাখলাম।}
\end{quote}

\begin{quote}
\emph{Āmi ruṭir jônyô āṭā mākhlām.}
\end{quote}

\begin{quote}
I kneaded flour for the rotis.
\end{quote}

Past (he/she …): \bn{তিনি} \bn{ময়দা} \bn{মাখলেন}, \emph{tini môydā mākhlen}.

\begin{cousinweb}
\textbf{\bn{মাখা}} (\emph{mākhā}) is a tadbhava: CDIAL lists Bengali \emph{mākhā}, \textquotedblleft{}to knead, to smear\textquotedblright{}, under Sanskrit \emph{mrakṣati} (CDIAL 10379).
\end{cousinweb}

\begin{grammarlens}[title={What you've built}]
One more word for describing what you lived and felt.
\end{grammarlens}

\subsection*{Your turn}
\begin{itemize}
\raggedright
  \item \emph{Say it:} \emph{mākhā}
  \item \emph{Say it:} \emph{mākhā}, once more
  \item \emph{Say it:} \emph{Āmi ruṭir jônyô āṭā mākhlām.} — I kneaded flour for the rotis.
  \item \emph{Recall:} say the Bengali for a peel, then the Bengali for to grate, then say \emph{mākhā} again
\end{itemize}

\begin{tcolorbox}[breakable,colback=teal!4,colframe=teal!35!black,title={Before you move on}]
What does \bn{মাখা} mean? (\textquotedblleft{}To knead; to mix by hand\textquotedblright{}.) Say it once more.
\end{tcolorbox}
\section[\texorpdfstring{\bn{কুচোনো} (kuchono)}{কুচোনো (kuchono)}]{\bn{কুচোনো} (kuchono) — to chop finely; to dice}

\label{lesson:BN-C513-kuchono}



\begin{quote}
Before the new one: say the Bengali for to grate, then the Bengali for to knead.
\end{quote}

\subsection*{You'll want to know: \bn{কুচোনো}}
\textbf{\bn{কুচোনো}} — \emph{kuchono} — \textquotedblleft{}to chop finely; to dice\textquotedblright{}.

\begin{quote}
\bn{আমি} \bn{তরকারির} \bn{জন্য} \bn{পেঁয়াজ} \bn{কুচোলাম।}
\end{quote}

\begin{quote}
\emph{Āmi tôrkārir jônyô pẽyāj kucholām.}
\end{quote}

\begin{quote}
I chopped the onions finely for the curry.
\end{quote}

Past (he/she …): \bn{তিনি} \bn{আদা} \bn{কুচোলেন}, \emph{tini ādā kucholen}.

\begin{cousinweb}
\textbf{\bn{কুচোনো}} (\emph{kuchono}) comes from a reconstructed \emph{kucyate}, \textquotedblleft{}contracts, shrinks\textquotedblright{}; CDIAL lists Bengali \emph{kucāna}, \textquotedblleft{}to shrivel, to cut into bits\textquotedblright{} (CDIAL 3221).
\end{cousinweb}

\begin{grammarlens}[title={What you've built}]
One more word for describing what you lived and felt.
\end{grammarlens}

\subsection*{Your turn}
\begin{itemize}
\raggedright
  \item \emph{Say it:} \emph{kuchono}
  \item \emph{Say it:} \emph{kuchono}, once more
  \item \emph{Say it:} \emph{Āmi tôrkārir jônyô pẽyāj kucholām.} — I chopped the onions finely for the curry.
  \item \emph{Recall:} say the Bengali for to grate, then the Bengali for to knead, then say \emph{kuchono} again
\end{itemize}

\begin{tcolorbox}[breakable,colback=teal!4,colframe=teal!35!black,title={Before you move on}]
What does \bn{কুচোনো} mean? (\textquotedblleft{}To chop finely; to dice\textquotedblright{}.) Say it once more.
\end{tcolorbox}
